(i) is the definition of the shadow. (ii) is Theorem \claimref{odd-extensions.doubling-or-complexification}(iii) with $e_{5}^{2} = +1$: the products square to $-e_{\mu}^{2}$, which is $(-1, +1, +1, +1)$, and generate a subalgebra of dimension sixteen inside the even subalgebra of dimension sixteen. (iii) The shadow's even part consists of even elements of $\Cl(2,3)$, so it lies in the intersection; the shadow's odd part consists of odd elements of $\Cl(2,3)$, since the generators $e_{1}, \ldots, e_{4}$ are odd there, so it meets the even subalgebra in zero; the intersection is exactly $\Cl^{0}(1,3)$, of dimension eight, and $\Cl^{0}(1,3) \cong \Cl(3,0)$ (\claimref{std-algebra.sta-matrix}) $\cong \Mat(2,\C)$ (\claimref{std-algebra.pauli-algebra}) is the Pauli algebra. For the spanning statement: $\omega$ times a blade of the shadow is, up to sign, the complementary blade of the shadow times $e_{5}$, so $\Cl(1,3) + \omega\,\Cl(1,3)$ contains every blade of $\Cl(2,3)$ (the script checks the dimension); this is the shadow's version of Theorem \claimref{odd-extensions.doubling-or-complexification}(ii), which decomposes by the even subalgebra instead. (iv) An automorphism carries a subalgebra to an isomorphic subalgebra, and $\Cl(1,3) \not\cong \Cl(3,1)$ by Fact \claimref{odd-extensions.imported-representatives-dimension-nearest-four-types}. With $\omega$ central and $\omega^{2} = -1$, $(\omega e_{\mu})^{2} = -e_{\mu}^{2}$ and the $\omega e_{\mu}$ anticommute pairwise, so they generate a $\Cl(3,1)$; this is the complex-linear identification of the two real forms and it is not an automorphism of the real algebra, since $\omega e_{\mu}$ is a $4$-blade, not a vector.

For $\Cl(1,4)$: Theorem \claimref{odd-extensions.doubling-or-complexification}(iii) with $e_{5}^{2} = -1$ gives $(e_{\mu}e_{5})^{2} = e_{\mu}^{2}$, and $\omega^{2} = +1$ gives $(\omega e_{\mu})^{2} = e_{\mu}^{2}$, so all three sets generate $\Cl(1,3)$. For the non-embedding: $\Cl(1,4) \cong A \oplus A$ with $A = \Cl(1,3) \cong \Mat(2,\HH)$ by Theorem \claimref{odd-extensions.doubling-or-complexification}(i) and Fact \claimref{odd-extensions.imported-representatives-dimension-nearest-four-types}. A unital homomorphism $\phi: \Mat(4,\R) \to A \oplus A$ composed with the projection to a summand on which $\phi(1) \neq 0$ is a unital homomorphism $\Mat(4,\R) \to A$; its kernel is a two-sided ideal of the simple algebra $\Mat(4,\R)$ not containing $1$, hence zero; so it is injective between real algebras of equal dimension sixteen, hence an isomorphism, contradicting $\Mat(4,\R) \not\cong \Mat(2,\HH)$ (Fact \claimref{odd-extensions.imported-representatives-dimension-nearest-four-types}). The statements for the mirror's extensions follow by exchanging the roles of $p$ and $q$ throughout. The script verifies the squares, the anticommutation, the dimensions of the generated subalgebras and of their intersections with the even subalgebra, and the Wick-rotated squares, for all four extensions.
